\documentclass[10pt,a4paper]{amsart}
\usepackage[margin=19mm]{geometry}
\usepackage{amsmath,amssymb,mathtools}
\usepackage[scr=rsfs,cal=euler]{mathalfa}
\usepackage{xfrac}
\usepackage[unicode,hidelinks,bookmarks=false]{hyperref}
\newcommand{\ie}{i.e.\ }
\newcommand{\cf}{cf.\ }
\newcommand{\resp}{respectively\ }
\newcommand{\Subsection}{\subsection}
\providecommand{\nobreakdash}{\nobreak\hbox{-}}
\setlength{\emergencystretch}{2em}
\multlinegap0pt
\usepackage{bm}
\usepackage{bbm}
\usepackage{dsfont}
\usepackage{xspace}
\usepackage{cancel}
\usepackage{mathtools}
\usepackage{amsthm}
\makeatletter
\@ifundefined{theorem}{\newtheorem{theorem}{Theorem}}{}
\@ifundefined{claim}{\newtheorem{claim}[theorem]{Claim}}{}
\@ifundefined{lemma}{\newtheorem{lemma}[theorem]{Lemma}}{}
\makeatother
\makeatletter
\newcommand{\Cov}{\operatorname{Cov}}
\def\Ga{{\Gamma}}
\def\La{{\Lambda}}
\newcommand{\bE}{\mathbb E}
\newcommand{\bR}{\mathbb R}
\newcommand{\bZ}{\mathbb Z}
\newcommand{\cE}{\mathcal{E}}
\newcommand{\cG}{\mathcal G}
\expandafter\ifx\csname manuallemma\endcsname\relax
\newenvironment{manuallemma}[1]{%
	\renewcommand\themanuallemmainner{#1}%
	\manuallemmainner
}{\endmanuallemmainner}
\fi
\expandafter\ifx\csname manualquestion\endcsname\relax
\newenvironment{manualquestion}[1]{%
	\renewcommand\themanualquestioninner{#1}%
	\manualquestioninner
}{\endmanualquestioninner}
\fi
\expandafter\ifx\csname manualtheorem\endcsname\relax
\newenvironment{manualtheorem}[1]{%
	\renewcommand\themanualtheoreminner{#1}%
	\manualtheoreminner
}{\endmanualtheoreminner}
\fi
\makeatother
\title[Correlation decay for U(1) lattice Higgs theory: the case of]{Correlation decay for U(1) lattice Higgs theory: the case of small mass}
\author{Sourav Chatterjee, Oren Yakir}
\date{Source: arXiv:2509.19176v1 (CC BY 4.0)}
\begin{document}
\maketitle

\noindent\emph{Source and licence.} Correlation decay for U(1) lattice Higgs theory: the case of small mass. Version arXiv:2509.19176v1 (CC BY 4.0), distributed under the Creative Commons Attribution 4.0 International licence, as recorded in the arXiv licence metadata for this exact version and shown on the abstract page https://arxiv.org/abs/2509.19176v1. Full rights notice: licenses/arxiv-2509.19176-CC-BY-4.0.txt.\par\smallskip

\noindent\textit{Editorial scope.} Counted unit: the Claim whose source label is \texttt{claim:bound\_of\_derivative\_of\_non\_gaussian\_correction\_with\_repsect\_to\_gaussian\_measure} in the article (source0.tex lines 1386--1396, source label \texttt{claim:bound\_of\_derivative\_of\_non\_gaussian\_correction\_with\_repsect\_to\_gaussian\_measure}), delivered together with the article's own complete original proof of it (source0.tex lines 1398--1432, one \texttt{proof} environment of 35 lines). No mathematics is written, repaired or re-derived: statement and proof are the article's own text line for line. Required context retained uncounted: eq:def\_of\_T\_beta (lines 409--412); claim:bound\_on\_function\_on\_collection\_of\_subsets (lines 846--858); lemma:gaussian\_integration\_by\_parts (lines 1206--1215); eq:def\_of\_D\_pi (lines 1212--1214); eq:def\_of\_mathcal\_X (lines 1319--1322); claim:estimate\_of\_derivative\_non\_gaussian\_correction\_X (lines 1324--1334); eq:derivative\_of\_Xi\_wrt\_Delta\_after\_leibnitz (lines 1381--1383). Every definition, display and label of those lines is retained unchanged. Omitted from this package and not redistributed: 1--408, 413--845, 859--1205, 1215--1318, 1323--1323, 1335--1380, 1384--1385, 1397--1397, 1433--1556. Nothing inside a retained excerpt is omitted or altered: every retained body is delivered exactly as the article's lines are written, so no omission receipt of any kind accompanies this package. Citations: the retained excerpts carry no citation command at all, so no bibliography entry of the article is transcribed and no reference is redistributed. Reference adaptation: none was needed and none was made; every reference inside the delivered unit resolves inside it, so no unresolved reference or citation is produced. Printed slips kept literal: no printed word, symbol, hypothesis or number of the counted statement or of its proof was altered. Layout and class: the article's own class is \texttt{article} with packages including amssymb, amsfonts, amsthm, amsmath, upref, booktabs, graphicx, verbatim, centernot, xcolor, multicol, enumitem, geometry, fancyhdr; the wrapper is the lane's pinned \texttt{amsart} editorial wrapper at 10pt on A4, which restates the theorem-like environments the retained text needs and adds the article's own macro definitions that the retained text needs (\texttt{\textbackslash Cov}, \texttt{\textbackslash Ga}, \texttt{\textbackslash La}, \texttt{\textbackslash bE}, \texttt{\textbackslash bR}, \texttt{\textbackslash bZ}, \texttt{\textbackslash cE}, \texttt{\textbackslash cG}), with extra wrapper packages (bm, bbm, dsfont, xspace, cancel, mathtools). The delivered numbering is the unit's own. Assets: no retained span contains a figure environment, a graphics-inclusion command, a PDF-page inclusion, a table or a TikZ picture; no image file is copied into this package, the package declares no asset dependency and no image file is redistributed with it. Licence: version arXiv:2509.19176v1 of this article is distributed under the Creative Commons Attribution 4.0 International licence, as recorded in the arXiv licence metadata for this exact version and shown on the abstract page https://arxiv.org/abs/2509.19176v1. A full rights notice is delivered at licenses/arxiv-2509.19176-CC-BY-4.0.txt. Characters: every retained excerpt is pure ASCII, so the delivered engine, XeLaTeX with its default Latin Modern text font, reports no missing character.
    \begin{equation}
        \label{eq:def_of_T_beta}
        T_\beta = \frac{\log^{d+2}(\beta)}{\sqrt{\beta}} \, ,
    \end{equation}

    \begin{claim}
    \label{claim:bound_on_function_on_collection_of_subsets}
        Let $f$ be a non-negative function on subsets of $L\bZ^d$ and set
        \[
        F(\mathbf{X}) = \prod_{j=1}^n f({\sf P}_j) \, ,
        \]
        where $\mathbf{X} = ({\sf P}_1,\ldots,{\sf P}_n)$ is a collection of such  
        subsets. For all $A\ge 1$ and for all $\widetilde{{\sf P}} \subset L\bZ^d$, we have
        \[
        \sum_{n\ge 1} \frac{1}{n!} \bigg(\sum_{\substack{\mathbf{X} = ({\sf P}_1,\ldots,{\sf P}_n) \\ {\sf P}_j \  \text{and $\widetilde{{\sf P}}$ touch}}} F(\mathbf{X}) \bigg) \le \frac{1}{A} \exp\Big(A \, |\widetilde{{\sf P}}| \cdot \|f\|_0\Big) \, .
        \]
        A similar bound also holds for functions defined on the faces of $L\bZ^d$. 
    \end{claim}

    \begin{lemma} \label{lemma:gaussian_integration_by_parts}
        Let $F:\bR^{E(\La)} \to \bR$ be a smooth function which grows, along with all of its finite partial derivatives, at most polynomial at infinity. Then
        \[
        \partial^\Ga \bE_{\mathbf{s}}^{\sf G}\big[ F(\boldsymbol{\theta}) \big] = \sum_{\Pi\in \mathcal{P}(\Ga)} \bE_{\mathbf{s}}^{\sf G}\Big[ \big(\prod_{\pi\in \Pi} D^\pi\big) F(\boldsymbol{\theta}) \Big] \, ,
        \]
        where
        \begin{equation} \label{eq:def_of_D_pi}
         D^\pi F = \frac{1}{2} \sum_{x,y\in \cE(\cG)} \big(\partial^\pi \Cov_{\mathbf{s}} (\theta_x,\theta_y) \big) \, \frac{\partial^2 F}{ \partial \theta_x \partial \theta_y} + \sum_{x\in \cE(\cG)} \big(\partial^\pi \, \bE_{\mathbf{s}}^{\sf G}[\theta_x]\big) \, \frac{\partial F}{\partial \theta_x} \, .
        \end{equation}
    \end{lemma}

    \begin{equation}
    \label{eq:def_of_mathcal_X}
        \mathcal{X} = X(\Delta)\cup (\mathcal{B}_1 \setminus \mathcal{B}) \, . 
    \end{equation}

    \begin{claim}\label{claim:estimate_of_derivative_non_gaussian_correction_X}
        Let $\Delta\subset P^\prime(\mathcal{G}_1)$ and let $\mathcal{X}$ be given by~\eqref{eq:def_of_mathcal_X}. Then for all $R \ge 1$ and for all $L$ large enough there exists $C_L>0$ so that  
        \[
        \bigg| \partial^\Delta \int_{\bR^{\cE(\mathcal{X})}} e^{-\beta V_{\Lambda,\mathcal{X}}(\boldsymbol{\theta};\mathbf{s}_\Delta)} \chi_{\mathcal{G}}(\boldsymbol{\theta}) \, {\rm d}\mu_{\mathcal{X},\eta,\mathbf{s}_\Delta}^{\sf G} \bigg| \le \nu(\beta) \, \exp\Big(-R|X(\Delta)| + C_L |\mathcal{B}_1|\Big)  \, , 
        \]
        where
        \[
        \lim_{\beta \to \infty} \nu(\beta) = 0
        \]
        uniformly in $R$ and $\Delta$. 
    \end{claim}

    \begin{multline} \label{eq:derivative_of_Xi_wrt_Delta_after_leibnitz}
    \partial^\Delta \int_{\bR^{\cE(\mathcal{X})}} e^{-\beta V_{\Lambda,\mathcal{X}}(\boldsymbol{\theta};\mathbf{s}_\Delta)} \chi_{\mathcal{G}}(\boldsymbol{\theta}) \, {\rm d}\mu_{\mathcal{X},\eta,\mathbf{s}_\Delta}^{\sf G} \\  = \sum_{\Upsilon\subset \Delta} \partial_{{\sf G}}^{\Upsilon}  \, \bE_{\mathcal{X},\eta,\mathbf{s}_{\Delta}}^{\sf G}\Big[ \Big( \prod_{{\sf F} \in \Delta\setminus \Upsilon} \sum_{\substack{p\in P(\La) \\ E(p)\cap {\sf F} \not= \emptyset }}\beta g({\sf d} \theta_p)\Big) \cdot e^{-\beta V_{\La,\mathcal{X}} (\boldsymbol{\theta}; \mathbf{s}_{\Delta})} \, \chi_{\mathcal{G}} (\boldsymbol{\theta}) \Big]
    \end{multline}

    \begin{claim} \label{claim:bound_of_derivative_of_non_gaussian_correction_with_repsect_to_gaussian_measure}
        For any $\Upsilon\subset \Delta$, $\Pi\in \mathcal{P}(\Upsilon)$ and for all $L$ large enough, we have
        \[
        \bigg| \,  \bE_{\mathcal{X},\eta,\mathbf{s}_{\Delta}}^{\sf G} \Big[ \big( \prod_{\pi\in \Pi} D^\pi\big) e^{-\beta V_{\La,\mathcal{X}} (\boldsymbol{\theta}; \mathbf{s}_{\Delta})} \, \chi_{\mathcal{G}} (\boldsymbol{\theta})  \Big] \bigg| \le C^{|\mathcal{X}|} \prod_{\pi\in \Pi} \big(\nu(\beta) \, (L^{d})^{|\pi|} \, e^{-c d(\pi)}\big)\, , 
        \]
        where $D^\pi$ is given by~\eqref{eq:def_of_D_pi}, $d(\pi)$ is the length of the shortest path in $\bZ^d$ which visits all faces of $\pi$ and 
        \[
        \lim_{\beta\to \infty} \nu(\beta) = 0 \, ,
        \]
        uniformly in $\Delta$. 
    \end{claim}

    \begin{proof}[Proof of Claim~\ref{claim:estimate_of_derivative_non_gaussian_correction_X}]
        Recall that the presence of the smooth cut-off $\chi_\mathcal{G}$ forces all $e\in \cE(\mathcal{G})$ to have $|\theta_e| \le 2 T_\beta$, where $T_\beta$ is given by~\eqref{eq:def_of_T_beta}. In particular, for plaquettes $p\in P^\prime(\mathcal{G}_1)$ we have 
        \[
        \beta |g({\sf d}\theta_p)| \lesssim \beta \,  (T_\beta)^4 \le \frac{\log^{4d+8}(\beta)}{\beta} \, .
        \]
        Plugging this bound into~\eqref{eq:derivative_of_Xi_wrt_Delta_after_leibnitz} and applying Lemma~\ref{lemma:gaussian_integration_by_parts} to deal with the $\partial_{\sf G}^{\Upsilon}$ derivatives, we arrive at 
        \begin{multline*}
            \bigg|\partial^\Delta \int_{\bR^{\cE(\mathcal{X})}} e^{-\beta V_{\Lambda,\mathcal{X}}(\boldsymbol{\theta};\mathbf{s}_\Delta)} \chi_{\mathcal{G}}(\boldsymbol{\theta}) \, {\rm d}\mu_{\mathcal{X},\eta,\mathbf{s}_\Delta}^{\sf G} \bigg| \\ \le C^{|\Delta|} \sup_{\Upsilon\subset \Delta} \bigg[ \Big(L^{d-1} \, \frac{\log^{4d+8}(\beta)}{\beta}\Big)^{|\Delta \setminus \Upsilon|} \sum_{\Pi\in \mathcal{P}(\Upsilon)} \bE_{\mathcal{X},\eta,\mathbf{s}_{\Delta}}^{\sf G} \Big[ \big( \prod_{\pi\in \Pi} D^\pi\big) e^{-\beta V_{\La,\mathcal{X}} (\boldsymbol{\theta}; \mathbf{s}_{\Delta})} \, \chi_{\mathcal{G}} (\boldsymbol{\theta})  \Big] \bigg] \, .
        \end{multline*}
        By Claim~\ref{claim:bound_of_derivative_of_non_gaussian_correction_with_repsect_to_gaussian_measure}, we can further bound the inner sum as
        \[
        \bigg|\sum_{\Pi\in \mathcal{P}(\Upsilon)} \bE_{\mathcal{X},\eta,\mathbf{s}_{\Delta}}^{\sf G} \Big[ \big( \prod_{\pi\in \Pi} D^\pi\big) e^{-\beta V_{\La,\mathcal{X}} (\boldsymbol{\theta}; \mathbf{s}_{\Delta})} \, \chi_{\mathcal{G}} (\boldsymbol{\theta})  \Big] \bigg| \le  C^{|\mathcal{X}|} \sum_{\Pi\in \mathcal{P}(\Upsilon)} \prod_{\pi\in \Pi} \Big(\nu(\beta) \, (L^d)^{|\pi|} \, e^{-c d(\pi)}\Big) \, .
        \]
        Since $|\mathcal{X}| \le |X(\Delta)| + |\mathcal{B}_1|$, we can combine both bounds and arrive at the inequality 
        \begin{multline} \label{eq:proof_of_claim:estimate_of_derivative_non_gaussian_correction_X_after_taking_derivatives}
             \bigg|\partial^\Delta \int_{\bR^{\cE(\mathcal{X})}} e^{-\beta V_{\Lambda,\mathcal{X}}(\boldsymbol{\theta};\mathbf{s}_\Delta)} \chi_{\mathcal{G}}(\boldsymbol{\theta}) \, {\rm d}\mu_{\mathcal{X},\eta,\mathbf{s}_\Delta}^{\sf G} \bigg| \\ \le C^{|X(\Delta)| + |\mathcal{B}_1|}  \sup_{\Upsilon\subset \Delta} \bigg[ \Big(L^{d-1} \, \frac{\log^{4d+8}(\beta)}{\beta}\Big)^{|\Delta \setminus \Upsilon|} \sum_{\Pi\in \mathcal{P}(\Upsilon)} \prod_{\pi\in \Pi} \Big(\nu(\beta) \, (L^d)^{|\pi|} \,  e^{-c d(\pi)}\Big) \bigg] \, .
         \end{multline}
         Now, for any $R \ge 1$ we have 
         \begin{multline*}
             \sum_{\Pi\in \mathcal{P}(\Upsilon)} \prod_{\pi\in \Pi} \Big(\nu(\beta) \, (L^d)^{|\pi|} \,  e^{-c d(\pi)}\Big) \\ = e^{-R|\Upsilon|} \sum_{\Pi\in \mathcal{P}(\Upsilon)} \prod_{\pi\in \Pi} \Big(\nu(\beta) \, e^{R|\pi|} \, (L^d)^{|\pi|} \,  e^{-c d(\pi)}\Big)  \\ \stackrel{\text{Claim}~\ref{claim:bound_on_function_on_collection_of_subsets}}{\le} e^{-R|\Upsilon|} \,  \nu(\beta) \,  \exp\Big(|\Upsilon| \cdot  \| (L^d)^{|\cdot|} \, e^{-cd(\cdot)} \|_{R}\Big) \, .
         \end{multline*}
         Since $d(\pi) \ge c L |\pi|$, we can choose $L$ large enough so that $\|(L^d)^{|\cdot|} \,  e^{-cd(\cdot)} \|_{R} \le 1$, and get that 
         \[
           \sum_{\Pi\in \mathcal{P}(\Upsilon)} \prod_{\pi\in \Pi} \Big(\nu(\beta) \, (L^d)^{|\pi|} \,  e^{-c d(\pi)}\Big)  \le e^{-(R-1)|\Upsilon|} \,  \nu(\beta) \, . 
         \]
         Plugging the above into~\eqref{eq:proof_of_claim:estimate_of_derivative_non_gaussian_correction_X_after_taking_derivatives} and taking $\beta$ so large such that 
         \[
         L^{d-1} \, \frac{\log^{4d+8}(\beta)}{\beta} \le e^{-2R} \, ,
         \]
         we obtain
         \[
         \bigg|\partial^\Delta \int_{\bR^{\cE(\mathcal{X})}} e^{-\beta V_{\Lambda,\mathcal{X}}(\boldsymbol{\theta};\mathbf{s}_\Delta)} \chi_{\mathcal{G}}(\boldsymbol{\theta}) \, {\rm d}\mu_{\mathcal{X},\eta,\mathbf{s}_\Delta}^{\sf G} \bigg|  \le \nu(\beta) \, C^{|X(\Delta)| + |\mathcal{B}_1|}   \, e^{-R|\Delta|} 
         \]
         from which the claim follows. 
    \end{proof}

\end{document}
